

\newcommand{\map}{\mathop{\mathit{map}}}
\newcommand{\reduce}{\mathop{\mathit{reduce}}}
\newcommand{\mapreduce}{\textsc{MapReduce}}
\newcommand{\shuffle}{\mathop{\mathit{shuffle}}}
\newcommand{\pig}{\textsc{Pig}}

\section{Programmer facilement en \mapreduce\,: \pig}


\subsection{Basics}
\begin{frame}
\frametitle{Pig Latin}

Un langage de haut niveau, type SQL, exécuté en  \mapreduce.

\vfill

``Requête'' \pig :  opération qui prend une \alert{relation} en entrée
et produit une autre relation en sortie.

\vfill

Les requêtes consistent en une \alert{séquence} d'instructions
\begin{redbullet}
 \item  \textit{LOAD} charge les données d'un fichier (HDFS)
       sous la forme d'une \alert{relation} (liste de n-uplets).
   \item Une série de transformations traite les données.
   \item L'instruction (optionnelle) \textit{STORE} écrit le résultat dans le système
       de fichiers (HDFS); on peut aussi faire un 
        \textit{DUMP} pour afficher à l'écran.
\end{redbullet}

\vfill

Les requêtes sont evaluées par une \alert{combinaison} de \textit{jobs} \mapreduce.
\end{frame}

\begin{frame}[fragile]
\frametitle{Tester \pig}

\pig est un programme Java qui s'exécute en ligne de commande.
\begin{redbullet}

\item Télécharger  depuis \url{http://pig.apache.org}.

\item Décompresser dans un répertoire, disons \texttt{pighome}

\item Placer le répertoire \texttt{pighome/bin} dans le chemin d'exécuton (\texttt{PATH})
\item Deux modes d'exécution
\begin{description}
\item[local] on lit les données sur le disque local, on évalue sans \mapreduce 
\item[MapReduce] on exécute dans un \textit{cluster} \mapreduce, en lisant depuis HDFS.
\end{description}



\end{redbullet}

Pour exécuter en local\,:
\begin{verbatim}
pig -x local
grunt>run <script>.pig
\end{verbatim}

\end{frame}

\begin{frame}[fragile]{Exemple de données en entrées (format tabulaire)}

Un fichier texte, un n-uplet par ligne, champs séparés par des tabulations.

\vfill

Fichier: \texttt{publishers.txt}

\footnotesize\begin{verbatim}
Fundations of Databases	Addison-Wesley	USA
Fundations of Databases	Vuibert	France
Web Data Management and Distribution	CUP	GB
\end{verbatim}

\vfill

Fichier: \texttt{books.txt}
\footnotesize\begin{verbatim}
1995	Fundations of Databases	Abiteboul
1995	Fundations of Databases	Hull
1995	Fundations of Databases	Vianu
2010	Web Data Management	Abiteboul
2010	Web Data Management	Manolescu
2010	Web Data Management	Rigaux
2010	Web Data Management	Rousset
2010	Web Data Management	Senellart
\end{verbatim}

\end{frame}

\begin{frame}[fragile]
\frametitle{Premier exemple: on regroupe les auteurs par livre.}



\begin{minted}{pig}
-- Chargement
books = load 'books.txt' 
    as (year: int, title: chararray, author: chararray) ;
-- Regroupement    
group_auth = group books by title;
-- Restruturation
authors = foreach group_auth generate group, books.author;
-- Affichage
dump authors;
\end{minted}

\begin{defblock}{Exécution}
\begin{verbatim}
pig -x local
grunt>run groupby.pig
\end{verbatim}
\end{defblock}

\end{frame}

\begin{frame}[fragile]
\frametitle{Modèle de données}

Modèle semi-structuré, type XML/JSON. Imbrication de structures ensemblistes (``bags'') et de n-uplets.

\vfill

Exemple: le \textit{bag} \texttt{group\_auth}.

\small{
\begin{minted}{pig}
(
  Fundations of Databases,
  { (1995,Fundations of Databases,Abiteboul),
    (1995,Fundations of Databases,Hull),
    (1995,Fundations of Databases,Vianu)
  }
)
\end{minted}
}

Imbrication sans limite, mais pas de références ou dépendance externe\,: un ``document'' autonome. 


\vfill

\alert{Objectif parallélisation}\,!
\end{frame}


\begin{frame}[fragile]
\frametitle{Un exemple complet: le fichier}

Fichier: \texttt{journal-small.txt}
\footnotesize\begin{verbatim}
2005    VLDB J. Model-based approximate querying in sensor networks.
1997    VLDB J. Dictionary-Based Order-Preserving String Compression.
2003	SIGMOD Record	Time management for new faculty.
2001    VLDB J. E-Services - Guest editorial.
2003	SIGMOD Record	Exposing undergraduate students to system internals.
1998    VLDB J. Integrating Reliable Memory in Databases.
1996    VLDB J. Query Processing and Optimization in Oracle Rdb
1996    VLDB J. A Complete Temporal Relational Algebra.
1994	SIGMOD Record	Data Modelling in the Large.
2002	SIGMOD Record	Data Mining: Concepts and Techniques - Book Review.
\end{verbatim}

\end{frame}


\begin{frame}[fragile]{Un exemple complet: le programme}

Regroupement des publications par année. 

\figSlideWithSize{pig-workflow}{10cm}

\vfill

\begin{minted}[frame=lines,
               baselinestretch=.9,
               fontsize=\footnotesize,
               framesep=2mm]{pig}
-- Chargement des documents de journal-small.txt
articles = load 'journal-small.txt' 
        as (year: chararray, journal:chararray, title: chararray) ;
sr_articles = filter articles by journal=='SIGMOD Record';
year_groups = group sr_articles by year;
count_by_year = foreach year_groups generate group, 
          count(sr_articles.title);
dump count_by_year;
\end{minted}

\vfill

À décomposer...
\end{frame}

\section{Les opérateurs \pig}

\begin{frame}[fragile]
\frametitle{Principaux opérateurs \pig}

\begin{small}
\begin{tabular}{lp{9cm}}
\hline
\textbf{Operator} & \textbf{Description} \\
\hline
\textit{foreach} & Applique une expression à chaque document\\ 
\textit{filter} &  Filtre les tuples sur certains critères\\ 
\textit{distinct} &  Elimination des doublons\\ 
\textit{join} &  Jointures de deux \textit{bags}\\ 
\textit{group} & Regroupement \\
\textit{cogroup} & Association de n-uplets provenant de deux sources\\ 
\textit{cross} &  Produit cartésien de deux sources\\ 
\textit{order} & Tri\\ 
\textit{limit} & Limite la taille des données traitées\\
\textit{union} &  Union de deux sources (dont les schémas peuvent être très différents)\\ 
\textit{split} & Partition d'une relation selon certains critères\\
\hline
\end{tabular}
\end{small}

\end{frame}

\begin{frame}[fragile]
\frametitle{Opérateur de ``dégroupage''}

\texttt{flatten} ``déploie'' une imbrication (\texttt{flatten.pig}).

\begin{minted}{pig}
-- Take the `authors' bag and flatten the nested set
flattened = foreach authors generate group, flatten($1);
\end{minted}

\vfill

Appliqué à la relation  \texttt{authors} créée précédemment, on obtient\,:


\small\begin{verbatim}
(Foundations of Databases,Abiteboul)
(Foundations of Databases,Hull)
(Foundations of Databases,Vianu)
(Web Data Management,Abiteboul)
(Web Data Management,Manolescu)
(Web Data Management,Rigaux)
(Web Data Management,Rousset)
(Web Data Management,Senellart)
\end{verbatim}
\end{frame}


\begin{frame}[fragile]
\frametitle{L'opérateur \texttt{cogroup}}

\vfill
Permet d'assembler des n-uplets sur un ensemble de critères.

\figSlideWithSize{pig-cogroup}{10cm}

\vfill
Exemple (\texttt{cogroup.pig})\,:

\begin{minted}{pig}
publishers = load 'webdam-publishers.txt' 
    as (title: chararray, publisher: chararray) ;
cogrouped = cogroup flattened by group,
                    publishers by title;
\end{minted}

\end{frame}


\begin{frame}[fragile]
\frametitle{Résultat (mis en forme)}
\vfill
Pour chaque valeur de regroupement, deux ensembles imbriqués,
provenant de chacune des sources.

\vfill

{\small\begin{verbatim}
(Foundations of Databases,
  { (Foundations of Databases,Abiteboul),
    (Foundations of Databases,Hull),
    (Foundations of Databases,Vianu)
  },
  {(Foundations of Databases,Addison-Wesley),
   (Foundations of Databases,Vuibert)
  }
)
\end{verbatim}
}

\vfill
Une sorte de jointure, dans un état intermédiaire, permis par le modèle semim-structuré.

\end{frame}


\begin{frame}[fragile]
\frametitle{Jointures}
\vfill
Come précédemment, mais produit un résultat ``mis à plat''. 

\vfill

On
effectue le produit cartésien des n-uplets imbriqués (\texttt{join.pig}).

\vfill

\begin{minted}{pig}
-- Take the 'flattened' bag, join with 'publishers'
joined = join flattened by group, publishers by title;
\end{minted}

\vfill

{\small\begin{verbatim}
(Foundations of Databases,Abiteboul,
      Fundations of Databases,Addison-Wesley)
(Foundations of Databases,Abiteboul,
      Fundations of Databases,Vuibert)
(Foundations of Databases,Hull,
      Fundations of Databases,Addison-Wesley)
(Foundations of Databases,Hull,
     ...
\end{verbatim}
}

\end{frame}


\begin{frame}[fragile]
\frametitle{Un exemple à tester}
\vfill
Avec ajout de filtres pour sélectionner certain n-uplets.


\vfill
               
\begin{minted}[frame=lines,
               baselinestretch=.9,
               fontsize=\footnotesize,
               framesep=2mm]{pig}
-- Chargement de books.txt
books = load 'books.txt' 
    as (year: int, title: chararray, author: chararray) ;
-- On prend les livres de Victor Vianu
vianu = filter books by author == 'Vianu';
--- Chargement de publishers.txt
publishers = load 'publishers.txt' 
    as (title: chararray, publisher: chararray) ;
-- Jointure sur le titre
joined = join vianu by title, publishers by title;
-- Groupement sur le nom de l'auteur
grouped = group joined by vianu::author;
-- Maintenant, comptons les editeurs (exercice: supprimer les doublons!)
count = foreach grouped generate group, 
           COUNT(joined.publisher);
\end{minted}

À vous de jouer.

\end{frame}

